\documentclass[../libro.tex]{subfiles}

\begin{document}

\ifSubfilesClassLoaded{\frontmatter}{}

% https://stackoverflow.com/a/51412802
\begingroup
\renewcommand{\thepage}{a}
\maketitle
\endgroup
\clearpage

\pagenumbering{alph}\setcounter{page}{2}%
\thispagestyle{empty}%

\begingroup
% \setlength{\parskip}{.25\baselineskip}
\setlength{\parindent}{1.5em}%for 12pt, or "xiaosi" (-4)

\esperanto{%
    \noindent%
    Ĉi tiu libro estas disponebla laŭ la \href{https://opensource.org/license/0bsd}{Nul-Kondiĉa Permesilo de BSD} (mallongigo de angle \textit{Berkeley Software Distribution}). Jen \href{https://github.com/septsea/permesiloj/blob/cxefa/0bsd.md}{\textbf{neoficiala} traduko de la permesilo en Esperanton}. Ĝi estas provizita por helpi pli multajn homojn kompreni la permesilon. Ĝi ne estis publikigita de la aŭtoro de la permesilo, kaj \textbf{ne} laŭleĝe deklaras la distribukondiĉojn por programaro, kiu uzas la permesilon—nur la originala angla teksto de la permesilo faras tion.%
}

\vspace{1.5ex}

\esperanto{%
    Aŭtorrajto © \lajaro{} \laauxtoro{}

    Permeso uzi, kopii, modifi kaj/aŭ distribui ĉi tiun programaron por iu ajn celo kun aŭ sen pago estas ĉi tie donita.

    \textbf{%
        La programaro estas provizita ``tiel, kiel estas'', kaj la aŭtoro ne donas iujn ajn garantiojn rilate al ĉi tiu programaro, inkluzive de ĉiuj implicitaj garantioj pri merkatkapableco kaj taŭgeco. En neniu okazo la aŭtoro estu respondeca pri iuj ajn specialaj, rektaj, nerektaj, aŭ konsekvencaj damaĝoj aŭ iaj ajn damaĝoj rezultantaj el perdo de uzado, datumoj aŭ profitoj, ĉu en ago de kontrakto, neglektado aŭ alia delikta ago, sekvantaj el aŭ rilataj al la uzado aŭ rendimento de ĉi tiu programaro.%
    }%
}

\vfill

\noindent%
This book is available under the \href{https://opensource.org/license/0bsd}{Zero-Clause BSD License}. Here is the licence.

\vspace{1.5ex}

Copyright © \lajaro{} \laauxtoro{}

Permission to use, copy, modify, and/or distribute this software for any purpose with or without fee is hereby granted.

\textbf{%
    The software is provided ``as is'' and the author disclaims all warranties with regard to this software including all implied warranties of merchantability and fitness. In no event shall the author be liable for any special, direct, indirect, or consequential damages or any damages whatsoever resulting from loss of use, data or profits, whether in an action of contract, negligence or other tortious action, arising out of or in connection with the use or performance of this software.%
}

\endgroup

\clearpage

\end{document}

% tex-fmt: off

I include the full text of the Zero-Clause BSD License (0BSD),
under which the book is available.
I also include an Esperanto translation
(for some reason that I have forgotten).

Here is a piece of trivia.
I have replaced "license"
(used as a common noun,
but not part of a proper noun, and not a verb)
by "licence".
I believe that it is not so good
to use "license" for both noun and verb forms.

Here is another piece of trivia.

0BSD is one of the so-called "public-domain-equivalent licences",
which grant public-domain-like rights and/or act as waivers,
and which are used to make copyrighted works
usable by anyone without conditions,
while avoiding the complexities of attribution
or licence compatibility
that occur with other licences.
Public-domain-equivalent licences exist
because some legal jurisdictions
do not provide for authors
to place their work in the public domain voluntarily,
but do allow them to grant arbitrarily broad rights
in the work to the public.
(I copied this paragraph from a free encyclopaedia.)

I admire Mr LI, Wenwei's work.
He said,
"I do not earn much by writing books:
the point is to help propagate the values of Open Knowledge
to the authors and publishers in China,
as well as to prove its profitability."
Not only does he write books,
he also releases the source code files of his published books.
He licenses his books under CC-BY-4.0
(in which "BY" means
that almost any use is permitted,
subject to providing credit and licence notice);
I go further to allow one "unlimited freedom" with my book
without requirements to include the copyright notice,
licence text,
or disclaimer in either source or binary forms
(or rather, "TeX or PDF files", in my case),
by using 0BSD.
(0BSD is a short and simple permissive licence
with no conditions, not even requiring preservation
of copyright and licence notices.
It gives one permissions to do
whatever one want with the code.)
(I did not use CC0, because CC0 is too long and too complex.)
I hope that using 0BSD facilitates the spread of my book.

Incidentally,
"Unlicense" and "MIT-0" are
also public-domain-equivalent licences.
For some reason that I have forgotten,
I translated these three licences into Esperanto:
https://github.com/septsea/permesiloj/
